% SPDX-License-Identifier: CC-BY-4.0
\documentclass[11pt,letterpaper]{article}

\usepackage[T1]{fontenc}
\usepackage{amsmath,amssymb}
\usepackage{tikz}
\usepackage{pgfplots}
\pgfplotsset{compat=1.18}
\usepackage{algorithm}
\usepackage{algpseudocode}
\usepackage{booktabs}
\usepackage{enumitem}
\usepackage[hidelinks]{hyperref}

\newcommand{\N}{\mathbb{N}}
\newcommand{\Z}{\mathbb{Z}}
\newcommand{\Q}{\mathbb{Q}}
\newcommand{\R}{\mathbb{R}}
\newcommand{\C}{\mathbb{C}}
\newcommand{\Map}[3]{#1 \colon #2 \to #3}
\newcommand{\Mod}[1]{#1\text{-}\mathrm{Mod}}
\newcommand{\GroupHom}{\mathrm{Hom}_{\mathbf{Grp}}}
\newcommand{\RingHom}{\mathrm{Hom}_{\mathbf{Ring}}}

\newtheorem{theorem}{Theorem}
\newtheorem{lemma}[theorem]{Lemma}
\newtheorem{corollary}[theorem]{Corollary}
\newtheorem{definition}[theorem]{Definition}
\newtheorem{remark}[theorem]{Remark}
\newenvironment{proof}{\noindent\textit{Proof.}}{\hfill\(\square\)}

\title{Residue criterion for the first affine hit}
\author{}
\date{\today}

\begin{document}
\maketitle

\begin{abstract}
Let \(q\), \(c\) and \(d\) be integers with \(|d|\ge 2\), and let
\(A_{j}(n)=q^{j}n+c\,U_{j}(q)\) be the affine state, where \(U_{j}(q)\) is the
sum of \(q^{0}\) through \(q^{j-1}\). The class of \(A_{j}(n)\) in
\(\Z/|d|\Z\) is the \(j\)-th iterate of \(r\mapsto qr+c\), starting at the
class of \(n\). Consequently \(d\) divides \(A_{j}(n)\) and divides no earlier
affine state if and only if that walk is zero at time \(j\) and nonzero at
every earlier time. Seeds in the same class have the same first hit.
\end{abstract}

\section{Affine states and the residue step}

Throughout, a parameter record consists of integers \(q\) and \(c\) and an
integer \(d\) with \(|d|\ge 2\). In particular \(d\neq 0\). Write \(m=|d|\) and
write \(\pi\) for the canonical ring homomorphism \(\Map{\pi}{\Z}{\Z/m\Z}\).
The symbol \(0\) in \(\Z/m\Z\) denotes the zero class. The offset sum is the
integer sequence defined by \(U_{0}(q)=0\) and
\[
U_{j+1}(q)=U_{j}(q)+q^{j}.
\]
Iteration of a self-map \(f\) starts at time zero: \(f^{0}(x)=x\) and
\(f^{j+1}(x)=f(f^{j}(x))\).

\begin{definition}\label{Def:AffineState}
The affine state of multiplier \(q\), offset \(c\), time \(j\in\N\) and seed
\(n\in\Z\) is
\[
A_{j}(q,c,n)=q^{j}n+c\,U_{j}(q).
\]
When the parameter record is fixed, write \(A_{j}(n)\) for \(A_{j}(q,c,n)\).
\end{definition}

\begin{definition}\label{Def:FirstHit}
The time \(j\) is a first hit of the seed \(n\) when \(d\) divides \(A_{j}(n)\)
and \(d\) divides \(A_{i}(n)\) for no \(i<j\).
\end{definition}

\begin{definition}\label{Def:ResidueStep}
The residue step of the parameter record is the self-map
\(\Map{\rho}{\Z/m\Z}{\Z/m\Z}\) given by
\[
\rho(r)=q r+c,
\]
where \(q\) and \(c\) are read in \(\Z/m\Z\).
\end{definition}

The raw step, which divides by \(d\) on the divisible class and applies
\(n\mapsto qn+c\) otherwise, is not used below. The first hit is a statement
about affine states only.

\begin{lemma}\label{Lem:TrapSumSeries}
For every integer \(q\) and every \(j\in\N\),
\[
(q-1)U_{j}(q)=q^{j}-1.
\]
\end{lemma}

\begin{proof}
Proceed by induction on \(j\). If \(j=0\), then \(U_{0}(q)=0\) and
\(q^{0}-1=0\), so both sides are zero. If the identity holds at \(j\), the
recurrence gives
\begin{align*}
(q-1)U_{j+1}(q)
&=(q-1)\bigl(U_{j}(q)+q^{j}\bigr)\\
&=(q-1)U_{j}(q)+(q-1)q^{j}\\
&=q^{j}-1+q^{j+1}-q^{j}\\
&=q^{j+1}-1.
\end{align*}
Thus the identity holds at \(j+1\).
\end{proof}

\begin{lemma}\label{Lem:TrapSumRec}
For every integer \(q\) and every \(j\in\N\),
\[
U_{j+1}(q)=q\,U_{j}(q)+1.
\]
\end{lemma}

\begin{proof}
Lemma~\ref{Lem:TrapSumSeries} rearranges to \(q^{j}=(q-1)U_{j}(q)+1\).
Substitute this expression into the recurrence:
\[
U_{j+1}(q)=U_{j}(q)+q^{j}
=U_{j}(q)+(q-1)U_{j}(q)+1
=q\,U_{j}(q)+1.
\]
\end{proof}

\begin{lemma}\label{Lem:AffineStateZero}
For all integers \(q\), \(c\) and \(n\), \(A_{0}(q,c,n)=n\).
\end{lemma}

\begin{proof}
By definition, \(A_{0}(q,c,n)=q^{0}n+c\,U_{0}(q)=n\), since \(q^{0}=1\) and
\(U_{0}(q)=0\).
\end{proof}

\begin{lemma}\label{Lem:AffineStateSucc}
For all integers \(q\), \(c\) and \(n\), and every \(j\in\N\),
\[
A_{j+1}(q,c,n)=q\,A_{j}(q,c,n)+c.
\]
\end{lemma}

\begin{proof}
Apply Lemma~\ref{Lem:TrapSumRec} and the definition:
\begin{align*}
A_{j+1}(q,c,n)
&=q^{j+1}n+c\,U_{j+1}(q)\\
&=q^{j+1}n+c\bigl(q\,U_{j}(q)+1\bigr)\\
&=q\bigl(q^{j}n+c\,U_{j}(q)\bigr)+c\\
&=q\,A_{j}(q,c,n)+c.
\end{align*}
\end{proof}

\section{Zero in the residue ring}

\begin{lemma}\label{Lem:CastZeroIffDvd}
Let \(d\) be an integer with \(m=|d|\ge 2\), and let \(a\in\Z\). Then
\(\pi(a)=0\) in \(\Z/m\Z\) if and only if \(d\) divides \(a\).
\end{lemma}

\begin{proof}
The kernel of \(\pi\) is the ideal \(m\Z\). Thus \(\pi(a)=0\) if and only if
\(m\) divides \(a\). It remains to identify the ideals \((m)\) and \((d)\).

Write \(d=\varepsilon m\) with \(\varepsilon\in\{1,-1\}\). If \(m\) divides \(a\),
then \(a=mk\) for some \(k\in\Z\), so \(a=d(\varepsilon k)\) and \(d\) divides
\(a\). If \(d\) divides \(a\), then \(a=dk\) for some \(k\in\Z\), so
\(a=m(\varepsilon k)\) and \(m\) divides \(a\). Therefore \(m\) divides \(a\) if
and only if \(d\) divides \(a\).
\end{proof}

\begin{theorem}\label{Thm:AffineCast}
Let \(q\), \(c\) and \(d\) be a parameter record. For every \(n\in\Z\) and every
\(j\in\N\),
\[
\pi\bigl(A_{j}(n)\bigr)=\rho^{j}\bigl(\pi(n)\bigr)
\]
in \(\Z/m\Z\), where \(m=|d|\).
\end{theorem}

\begin{proof}
Fix the parameter record and the seed \(n\), and induct on \(j\).

If \(j=0\), Lemma~\ref{Lem:AffineStateZero} gives \(A_{0}(n)=n\). The zeroth
iterate is the identity, so both sides equal \(\pi(n)\).

Suppose the identity holds at time \(j\). Lemma~\ref{Lem:AffineStateSucc} gives
\(A_{j+1}(n)=q\,A_{j}(n)+c\). Since \(\pi\) preserves addition and
multiplication,
\[
\pi\bigl(A_{j+1}(n)\bigr)
=\pi\bigl(q\,A_{j}(n)+c\bigr)
=q\,\pi\bigl(A_{j}(n)\bigr)+\pi(c).
\]
The right-hand side is \(\rho\) applied to \(\pi(A_{j}(n))\). The inductive
hypothesis replaces that class by \(\rho^{j}(\pi(n))\), and one further
application of \(\rho\) is the iterate of length \(j+1\).
\end{proof}

\begin{lemma}\label{Lem:AffineDvdIffResidueZero}
Let \(q\), \(c\) and \(d\) be a parameter record. For every \(n\in\Z\) and every
\(j\in\N\), \(d\) divides \(A_{j}(n)\) if and only if
\(\rho^{j}(\pi(n))=0\) in \(\Z/|d|\Z\).
\end{lemma}

\begin{proof}
By Lemma~\ref{Lem:CastZeroIffDvd}, \(d\) divides \(A_{j}(n)\) if and only if
\(\pi(A_{j}(n))=0\). Theorem~\ref{Thm:AffineCast} replaces
\(\pi(A_{j}(n))\) by \(\rho^{j}(\pi(n))\).
\end{proof}

\section{First hits}

\begin{theorem}\label{Thm:FirstHitIffResidue}
Let \(q\), \(c\) and \(d\) be a parameter record. For every \(n\in\Z\) and every
\(j\in\N\), the time \(j\) is a first hit of \(n\) if and only if
\(\rho^{j}(\pi(n))=0\) and \(\rho^{i}(\pi(n))\neq 0\) for every \(i<j\).
\end{theorem}

\begin{proof}
Write \(m=|d|\). By Lemma~\ref{Lem:AffineDvdIffResidueZero}, for every
\(i\in\N\) the integer \(d\) divides \(A_{i}(n)\) if and only if
\(\rho^{i}(\pi(n))=0\).

Suppose \(j\) is a first hit. Then \(d\) divides \(A_{j}(n)\), so
\(\rho^{j}(\pi(n))=0\). If some \(i<j\) satisfied \(\rho^{i}(\pi(n))=0\), the
same equivalence would give \(d\mid A_{i}(n)\), contradicting the first-hit
assumption. Thus \(\rho^{i}(\pi(n))\neq 0\) for every \(i<j\).

Conversely, suppose \(\rho^{j}(\pi(n))=0\) and \(\rho^{i}(\pi(n))\neq 0\) for
every \(i<j\). The equivalence gives \(d\mid A_{j}(n)\). If \(d\) divided
\(A_{i}(n)\) for some \(i<j\), the equivalence would give
\(\rho^{i}(\pi(n))=0\), which is excluded. Therefore \(j\) is a first hit.
\end{proof}

\begin{lemma}\label{Lem:ResidueIterateCongruent}
Let \(q\), \(c\) and \(d\) be a parameter record, and let \(n,n'\in\Z\) satisfy
\(\pi(n)=\pi(n')\) in \(\Z/|d|\Z\). Then for every \(i\in\N\),
\[
\rho^{i}\bigl(\pi(n)\bigr)=\rho^{i}\bigl(\pi(n')\bigr).
\]
\end{lemma}

\begin{proof}
Induct on \(i\). At \(i=0\) both sides are the common class \(\pi(n)\). If the
iterates agree at time \(i\), apply \(\rho\) to both sides. The resulting
classes are the iterates of length \(i+1\).
\end{proof}

\begin{lemma}\label{Lem:FirstHitOfCongruent}
Let \(q\), \(c\) and \(d\) be a parameter record, and let \(n,n'\in\Z\) satisfy
\(\pi(n)=\pi(n')\). If \(j\) is a first hit of \(n\), then \(j\) is a first hit
of \(n'\).
\end{lemma}

\begin{proof}
Theorem~\ref{Thm:FirstHitIffResidue} converts the first hit of \(n\) into the
statement that \(\rho^{j}(\pi(n))=0\) and \(\rho^{i}(\pi(n))\neq 0\) for every
\(i<j\). Lemma~\ref{Lem:ResidueIterateCongruent} replaces every one of these
classes by the corresponding iterate of \(\pi(n')\).
Theorem~\ref{Thm:FirstHitIffResidue}, read in the opposite direction for the
seed \(n'\), returns that \(j\) is a first hit of \(n'\).
\end{proof}

\begin{theorem}\label{Thm:FirstHitCongruent}
Let \(q\), \(c\) and \(d\) be a parameter record. If \(n,n'\in\Z\) satisfy
\(\pi(n)=\pi(n')\) in \(\Z/|d|\Z\), then for every \(j\in\N\) the time \(j\) is
a first hit of \(n\) if and only if it is a first hit of \(n'\).
\end{theorem}

\begin{proof}
If \(j\) is a first hit of \(n\), Lemma~\ref{Lem:FirstHitOfCongruent} gives the
same conclusion for \(n'\). If \(j\) is a first hit of \(n'\), the same lemma
applies with the seeds exchanged: the hypothesis \(\pi(n)=\pi(n')\) is
symmetric, so \(j\) is a first hit of \(n\).
\end{proof}

\begin{remark}
The three theorems do not assert that a first hit exists. A residue walk on
the finite ring \(\Z/m\Z\) may avoid zero at every time; the criterion then
says that no affine state is divisible by \(d\). They also do not identify
\(A_{j}(n)\) with a raw iterate, and they say nothing about the raw orbit at
any time after a division. No bound on a hitting time, no return exponent, and
no density or convergence statement is proved.
\end{remark}

\end{document}
